\documentclass[sigconf,anonymous,review]{acmart}
\setcopyright{none}
\setlength{\emergencystretch}{2em}
\settopmatter{printacmref=false}
\renewcommand\footnotetextcopyrightpermission[1]{}
\acmConference[RAIOps4Fin '26]{Workshop on Applied Responsible AI Operations for Finance}{November 2026}{Milan, Italy}
\acmYear{2026}
\title{Validating Portfolio Rewards: Preference Drift, Optimisation Error and Historical Performance}
\author{Anonymous Author(s)}
\newcommand{\E}{\mathbb{E}}
\newcommand{\Var}{\operatorname{Var}}
\newcommand{\Cov}{\operatorname{Cov}}
\newcommand{\eq}{\mathrm{eq}}
\newtheorem{proposition}{Proposition}
\newtheorem{theorem}{Theorem}
\begin{document}
\begin{abstract}
A portfolio model can pass an optimisation check while implementing a different risk preference from the one it was designed to follow. We develop a validation procedure that compares the financial reference, the exact reward optimum and the finite-budget policy on the same pre-trade states. For constrained mean--variance allocation, an auxiliary-variable identity shows that calibrating a squared-return penalty also changes the effective turnover cost. Joint calibration restores the fitted reference; risk-only calibration generally does not. We implement the procedure on ten US industry portfolios from 1990 to 2025, with 432 monthly decisions and 17 robustness settings. The audit distinguishes three findings: fitted-policy recovery is numerically precise, finite-budget optimisation error can exceed reward-design error, and realised performance gains remain statistically inconclusive. The resulting records make preference fidelity and solver accuracy inspectable without treating either as evidence of investment superiority.
\end{abstract}
\ccsdesc[500]{Computing methodologies~Reinforcement learning}
\ccsdesc[300]{Applied computing~Economics}
\keywords{portfolio allocation, reward alignment, mean--variance, reinforcement learning, transaction costs}
\maketitle

\section{Introduction}
A validation report for a learned allocator should answer a simple question: what preference does the deployed reward implement? Return, variance, squared return and turnover penalties can look similar in a specification while inducing different trades. A successful optimiser establishes only that the specified reward has been optimised. A favourable backtest answers a different question again.

We turn this distinction into a reproducible model-validation procedure. The procedure stores the intended objective, the reward-induced allocation, numerical optimality certificates and subsequent realised returns. Comparisons use common pre-trade states so that a reward change is not confused with differences accumulated along two trading paths.

Our case study is a constrained mean--variance allocator with proportional turnover costs. The mean auxiliary in MVPI\ \cite{zhang2021} gives an exact identity: adjusting the quadratic risk coefficient changes the effective cost preference unless the turnover coefficient is adjusted too. This supplies a concrete test case with a known fitted reference.

The historical evidence demonstrates both the use and the limit of the procedure. Joint calibration matches the reference across 17 settings, but the main realised performance intervals include zero. The contribution is an implemented validation method with independently checked records, rather than a claim of production deployment or regulatory certification.

\section{Portfolio objectives and validation}
Direct utility optimisation is established in financial learning. Simon et al.\ \cite{simon2026}, Section 4.1, Eq.\ (17), use centred mean--variance preferences, with turnover costs in Section 3.6. Their deep portfolio policy and US stock-characteristic data differ from the common industry-portfolio environment used here. Cong et al.\ \cite{cong2026}, Eqs.\ (19)--(20), optimise trajectory objectives through pathwise gradients. These approaches show that objective-aware learning is feasible.

Other rewards encode different preferences. Jiang et al.\ \cite{jiang2025}, Eq.\ (8), use net log return minus a variance penalty; Eq.\ (19) adds trading restrictions. We adapt that preference using monthly returns, hard simplex constraints and a first-order turnover charge. The comparison changes the reward while holding the estimated state fixed; it does not reproduce their CNN/WaveNet DPG system. Their earlier TD3 study\ \cite{jiang2024}, Eq.\ (4), combines portfolio value, variance and trading costs, motivating explicit treatment of all three terms.

The calibration below follows the mean auxiliary in MVPI\ \cite{zhang2021}. Known-policy validation also has precedents: Chaouki et al.\ \cite{chaouki2020}, Section 3, compare learned and analytical controls, while Gollart and Okhrin\ \cite{gollart2025}, Section 4.2.1 and Table 1, align terminal utility with a DDPG problem. Our contribution is a reproducible test of risk--cost preference drift and a comparison of its magnitude with optimisation error and historical uncertainty.

\section{Constrained alignment with turnover costs}
Let $m$ and $\Sigma\succ0$ be estimated excess-return moments available before trading, $M=\Sigma+mm^\top$, and $\nu=m^\top M^{-1}m$. The feasible set $C$ is the simplex, $p$ is the pre-trade allocation, and $c(u)=\|u-p\|_1$. For risk aversion $\varphi>0$ and cost coefficient $\lambda\ge0$, define
\begin{equation}\label{eq:target}
 F(u)=m^\top u-\varphi u^\top\Sigma u-\lambda c(u),\quad u\in C.
\end{equation}
Its unique optimum $u^*$ is a \emph{fitted reference}, not the unknown true market optimum. The reference specifies a one-period decision under estimated moments.

\begin{proposition}\label{prop:joint}
Let $y^*=m^\top u^*$ and $A=1+2\varphi y^*>0$. The second-moment reward objective
\begin{equation}
 Q(u)=m^\top u-\frac{\varphi}{A}u^\top M u-\frac{\lambda}{A}c(u)
\end{equation}
has the same unique maximiser as $F$.
\end{proposition}
\begin{proof}
For all feasible $u$,
\begin{equation}\label{eq:identity}
 AQ(u)=F(u)-\varphi(m^\top u-y^*)^2+\varphi(y^*)^2.
\end{equation}
At $u^*$ the square vanishes and $F$ is maximal. Since $A>0$, the assertion follows.
\end{proof}
The argument applies to any feasible set with a global maximiser. Convexity and positive definite covariance ensure uniqueness here. The identity follows the mean auxiliary used in MVPI\ \cite{zhang2021}.

If only the risk coefficient is divided by $A$, while the cost coefficient stays $\lambda$, the right-hand side of Eq.\ \eqref{eq:identity} gains the extra term
\begin{equation}\label{eq:costdistortion}
 -2\varphi y^*\lambda c(u).
\end{equation}
When $y^*>0$, this increases the effective cost preference. It need not imply monotone realised turnover on the strategy's evolving path. The unconstrained formula $\varphi(1-\nu)$ likewise need not calibrate a constrained optimum.

Joint calibration uses $u^*$ and therefore serves as a certificate of alignment with a known reference. As a constructive prior-art comparator, we also alternate maximising
\begin{equation}
 (1+2\varphi y)m^\top u-\varphi u^\top Mu-\lambda c(u)
\end{equation}
and setting $y=m^\top u$. This is an adaptation of the MVPI mean auxiliary to the constrained one-period problem. Including the term $-\varphi y^2$ gives a jointly strictly concave objective, since its quadratic block has Schur complement $\Sigma\succ0$. Exact block updates converge to the same target. This provides a constructive comparator with the same optimum.

\section{Historical study}
\subsection{Data and protocol}
We use Kenneth French's ten value-weighted US industry portfolio monthly total returns and the Fama--French risk-free series\ \cite{french}. The archived downloads have a recorded SHA-256 manifest. Estimation data start January 1970; evaluation covers January 1990--December 2025 (432 months). This is a retrospective study on the August 2026 revised vintage, not a point-in-time constituent dataset. Industry portfolios are research portfolios, not directly tradeable funds; their internal reconstruction costs are not included. The provider's CRSP processing and risk-free source changes also form part of this data vintage.

For month $t$, use only the previous 120 months of excess returns to estimate $m$ and the unbiased covariance, shrunk 10\% toward its diagonal. All strategies are long-only and fully invested. The primary settings are $\varphi=5$ and $\lambda=0.001$ (10 basis points per unit of L1 turnover). Pre-trade weights drift with realised gross asset returns. Net portfolio return uses the explicit first-order convention $u_t^\top R_t-\lambda\|u_t-p_t\|_1$, not a market-impact model or exact intramonth execution equation. The experimental protocol fixed these choices before strategy outcomes were inspected.

\subsection{Comparators and measurements}
We solve the covariance-MV target, raw second moment ($\kappa=\varphi$), unconstrained scalar calibration ($\kappa=\varphi(1-\nu)$), reference-based risk-only calibration, joint calibration, and MVPI adaptation. Equal weight and minimum variance provide conventional portfolio controls. Log utility and log-minus-variance adapt the reward preference in Jiang et al.\ \cite{jiang2025}; every method uses the same monthly data and feasible set. The comparison tests induced preferences under common estimates, not which published learning architecture is best.

For alignment, all methods face the target strategy's \emph{same pre-trade state}. We report Euclidean allocation distance and fitted regret $F(u^*)-F(u)$. Each strategy also follows its own pre-trade states for realised performance. This avoids conflating state divergence with the local reward comparison. Quadratic solutions have independent linear-programme first-order certificates; the largest saved common-state certificate is $4.27\times10^{-9}$ in objective units. The observed agreement of calibrated policies is considerably tighter than the acceptance tolerance.

Realised metrics are monthly excess-return MV, annualised excess Sharpe, return, volatility, turnover and drawdown. Uncertainty for paired MV differences uses 1,000 circular moving-block bootstrap resamples, blocks of 12 months, seed 8128. These intervals describe the observed strategy-return sequences, not full estimator-refitting or model-selection uncertainty. Robustness spans 17 prespecified combinations: $\varphi\in\{2,5,10\}$ and cost $\{0,10,25\}$ bp at the primary estimator, plus windows $\{60,120,240\}$ and covariance shrinkage $\{0,0.1,0.5\}$ at the primary preferences. The overlapping primary case is counted once.

\begin{table*}[t]
\centering
\caption{Reward-formulation comparison on historical data, 1990--2025. Published rewards are adapted to a common environment; complete published algorithms are not reproduced. Bold identifies our joint-calibration construction, not the best realised performance. Distance and fitted regret use common states; realised metrics use each strategy's own path. Gap and MV score are monthly basis-point units. Calibration improves model alignment; the table does not establish superior investment performance.}
\label{tab:historical}
\begin{tabular}{lrrrr}\toprule
Method & Allocation distance & Fitted MV gap (bp) & Excess Sharpe & Realised MV (bp)\\\midrule
Target & 0.00000 & 0.0000 & 0.570 & -8.497 \\
Raw second moment & 0.00843 & 0.0148 & 0.573 & -7.182 \\
Unconstrained scalar & 0.00540 & 0.0043 & 0.568 & -8.882 \\
Risk-only calibration & 0.00510 & 0.0033 & 0.571 & -8.365 \\
\textbf{Joint calibration (ours)} & \textbf{0.00000} & \textbf{0.0000} & \textbf{0.570} & \textbf{-8.497} \\
MVPI (Zhang et al., 2021)~\cite{zhang2021} & 0.00000 & 0.0000 & 0.570 & -8.497 \\
Log--variance (Jiang et al., 2025)~\cite{jiang2025} & 0.01286 & 0.0287 & 0.572 & -7.123 \\
Log utility & 0.54848 & 39.7855 & 0.344 & -133.295 \\
Equal weight & 0.43495 & 46.2626 & 0.622 & -11.893 \\
Minimum variance & 0.23843 & 9.0252 & 0.645 & 5.274 \\
\bottomrule\end{tabular}

\end{table*}

\subsection{Results}
Table~\ref{tab:historical} shows that raw second-moment optimisation has mean allocation distance 0.00843. Unconstrained scalar calibration reduces this to 0.00540; reference-based risk-only calibration retains 0.00510. Joint calibration and MVPI agree with the fitted reference to numerical precision. The adapted log-minus-variance objective has distance 0.01286; log utility alone is much farther away. The differences measure departure from the specified MV preference; log utility represents a different investment objective.

Across all 17 settings, the largest mean joint-calibration distance is $1.54\times10^{-7}$; raw mean distances range from 0.00157 to 0.03682. Thus constraints and trading costs do not break the joint identity, while unconstrained scalar calibration is not a universal repair. Figure~\ref{fig:historical} relates alignment to historical outcomes, and Figure~\ref{fig:robust} gives robustness and uncertainty.

Joint calibration has realised monthly MV lower than raw by 1.315 bp, with 95\% interval $[-3.318,0.343]$ bp. Against scalar calibration the difference is 0.385 bp. Its interval is $[-0.570,1.351]$ bp. Against risk-only calibration the difference is $-0.131$ bp, with interval $[-0.889,0.634]$ bp. All include zero. Minimum variance has both a better realised MV score and higher Sharpe than the fitted target in the primary sample. This illustrates estimation risk: precise implementation of estimated preferences cannot guarantee the best realised allocation. Equal weight has higher Sharpe than the target but lower realised MV, so Sharpe ranking alone does not answer even the ex-post MV comparison.

\begin{figure*}[p]
\centering\includegraphics[width=.94\textwidth]{figures/historical.pdf}
\caption{Left: mean common-state allocation distance; the logarithmic axis exposes numerical agreement for joint calibration and MVPI. Right: realised Sharpe and monthly MV need not rank policies alike. Equal weight has higher Sharpe but lower MV than the fitted target.}
\Description{Allocation distances on the left and realised Sharpe versus mean-variance on the right.}
\label{fig:historical}
\end{figure*}
\begin{figure*}[p]
\centering\includegraphics[width=.94\textwidth]{figures/robustness.pdf}
\caption{Left: alignment across the 17 prespecified settings, ordered by window, risk penalty, cost and shrinkage. Right: paired 12-month-block bootstrap intervals all include zero. Exact fitted alignment does not imply a statistically clear realised MV advantage.}
\Description{Seventeen-setting allocation-distance curves and three paired confidence intervals crossing zero.}
\label{fig:robust}
\end{figure*}

\subsection{Separating optimisation from reward design}
A finite-budget experiment assigns independent softmax policy logits to each of the 432 historical common states. Adam uses learning rate 0.05, initial logits $N(0,0.1^2)$, five matched seeds and 50, 250, 1,000 and 5,000 gradient steps. Each objective uses only its state's lagged estimates. Independent logits isolate optimisation error on each fitted state. They do not test generalisation across states, as a predictive neural policy would.

Let $u_r$ maximise the specified reward and $\widehat u_r$ be the finite-budget allocation. We report the nonnegative own-reward optimisation gap separately from
\begin{align}
 F(u^*)-F(\widehat u_r)
 &= [F(u^*)-F(u_r)]\nonumber\\
 &\quad+[F(u_r)-F(\widehat u_r)].
\end{align}
The second bracket is a \emph{signed} target-value difference; it is not the own-reward optimisation gap. At 5,000 steps, mean own-reward optimisation gaps remain approximately $4\times10^{-6}$, exceeding the exact raw design gap of $1.48\times10^{-6}$. Thus finite-budget error can dominate the reward discrepancy. Figure~\ref{fig:learning} shows why these must be kept separate: improving the raw-reward optimiser does not remove its exact design gap. The target and joint reward have the same limiting optimum but can exhibit different finite-budget trajectories because their objective geometries differ.

\begin{figure*}[p]
\centering\includegraphics[width=.94\textwidth]{figures/learning.pdf}
\caption{Finite-budget softmax optimisation on historical fitted states. Lines are five-seed means, shading the seed range. Left uses each method's own reward; right uses the common MV target. The dotted line is the exact raw reward-design gap. This diagnostic measures optimisation accuracy, not generalisation.}
\Description{Own-reward and target gaps decline with gradient steps; the raw target gap remains above its design floor.}
\label{fig:learning}
\end{figure*}

\section{A reproducible validation record}
For each decision date, the record stores the estimation window, preference coefficients, pre-trade allocation, chosen weights, realised return and numerical certificate. The raw data snapshots have recorded hashes. Reconstructing the common state and cost calculation from these records provides a check independent of the strategy runner.

Three comparisons have distinct interpretations. A nonzero gap between the reward optimum and the fitted reference identifies preference drift. A gap between a finite-budget policy and its own reward optimum identifies optimisation error. A difference in subsequent realised returns measures investment outcomes under estimation uncertainty. Combining these quantities into one score would obscure which part of the model requires attention.

The numerical tolerance is an implementation choice, not a business acceptance threshold. Quadratic iterates are accepted only after a feasible-set first-order certificate is at most $2\times10^{-7}$; the largest saved common-state certificate in the primary run is $4.27\times10^{-9}$. Risk-only calibration has mean fitted regret $3.35\times10^{-7}$, so the recorded numerical precision matters when interpreting such small effects. A model owner must separately decide whether the allocation discrepancy is economically material.

The implementation checks all 4,320 primary strategy decisions for lagged estimation dates, simplex feasibility, drift, cost accounting and common-state regret. The research test suite contains 31 passing tests, alongside ten original verification checks. These checks support the reported calculations; they do not establish live execution quality or a complete operational control system.
\section{Reconstructing the historical experiment}
\subsection{Information available at a decision}
Write $R_t\in\mathbb R^{10}$ for the vector of monthly total returns and $r_{f,t}$ for the risk-free return. For a window of length $L$, the information used at month $t$ is $\{R_s,r_{f,s}:t-L\le s<t\}$. In particular, neither $R_t$ nor a full-sample covariance enters the allocation. Let $X_s=R_s-r_{f,s}\mathbf1$. The sample mean and covariance are
\begin{align}
 m_t&=L^{-1}\sum_{s=t-L}^{t-1}X_s,\\
 S_t&=(L-1)^{-1}\sum_{s=t-L}^{t-1}(X_s-m_t)(X_s-m_t)^\top.
\end{align}
The implemented covariance is $\Sigma_t=(1-h)S_t+h\operatorname{diag}(S_t)$, with $L=120$ and $h=0.1$ in the primary experiment. The mean is not shrunk. Thus every reward sees identical estimates; differences are not attributable to different forecasting models. Positive definiteness is checked before optimisation. Historical dependence can be present in the observed returns without changing the algebra on a fixed estimated state. It does change how much confidence one should place in that estimated state as a forecast.

The raw files contain percent returns, which are divided by 100 before any objective is evaluated. Monthly records are aligned by date; missing-value sentinels, duplicate dates and missing joined observations are rejected. Estimation data begin in 1970 to support the longest 240-month window at the first 1990 decision. Observations after December 2025 are excluded even though the retrieved files extend further. The saved date range therefore fixes the comparison independently of the most recent available month.

\subsection{Two uses of the same historical month}
The alignment calculation and the investment calculation answer different questions. For alignment, all rewards receive the reference strategy's pre-trade portfolio $p_t^*$. Define $u_t^*$ as the fitted target allocation and $u_t^r$ as the allocation induced by reward $r$ on that common state. We record
\begin{equation}
 d_t^r=\|u_t^r-u_t^*\|_2,\qquad
 g_t^r=F_t(u_t^*)-F_t(u_t^r).
\end{equation}
The first measures the change in portfolio weights, while the second measures the economic cost under the stated fitted preference. A small objective gap can accompany a visible weight difference when the objective is locally flat. Neither quantity uses the month's subsequent realised return.

For investment performance, each method instead carries its own pre-trade portfolio $p_t^r$. It solves its reward again on that state, giving $\bar u_t^r$. The recorded turnover and net return are
\begin{equation}
 T_t^r=\|\bar u_t^r-p_t^r\|_1,\qquad
 r_t^r=(\bar u_t^r)^\top R_t-\lambda T_t^r.
\end{equation}
All methods start from equal weights. Before the next decision, holdings drift according to
\begin{equation}
 p_{t+1,i}^r=
 \frac{\bar u_{t,i}^r(1+R_{t,i})}{1+(\bar u_t^r)^\top R_t}.
\end{equation}
This uses gross asset returns to update relative holdings and subtracts costs separately in the first-order portfolio-return approximation. Charging exact self-financing costs would require a different wealth-accounting model. Using each method's own holdings in the alignment calculation would mix reward disagreement with the consequences of earlier decisions; the common-state record avoids that confounding.

\subsection{What the published baselines represent}
The MVPI comparator implements the mean auxiliary of Zhang et al.\ \cite{zhang2021} in the constrained one-period setting. Starting from $y=0$, it alternates a portfolio maximisation with the update $y\leftarrow m^\top u$. The portfolio step retains the original turnover coefficient. The procedure stops when the change in $y$ is below $10^{-10}$, with a maximum of 200 iterations. It computes the reference without first inserting the reference allocation into the reward. Joint calibration, by contrast, uses the already computed reference mean to demonstrate exact alignment. Their numerical agreement is a consistency check between two constructions, not an advantage over MVPI.

The Jiang et al.\ \cite{jiang2025} adaptation maximises
\begin{equation}
 \frac1L\sum_{s=t-L}^{t-1}\log(1+u^\top R_s)
 -\varphi u^\top\Sigma_tu-\lambda\|u-p_t\|_1.
\end{equation}
Setting its variance coefficient to zero gives the log-utility comparator. Logarithmic utility uses total returns, while the fitted MV objective uses excess returns. Both formulations receive the same historical window and feasible set. This is a deliberate comparison of preferences: logarithmic utility is not an algebraic approximation to the specified MV target. Hard simplex constraints and the linear turnover charge replace the source paper's soft restrictions and log cost factor. Its published returns are not pooled with ours, and its predictive architecture is not evaluated.

Equal weight supplies a fixed-allocation control. The minimum-variance control removes the expected-return term from the fitted target while retaining the risk scaling and turnover charge. It therefore tests whether using the estimated mean helps realised performance under this protocol. Its strong primary outcome is evidence that estimation quality matters even when reward alignment is exact.

\subsection{Numerical and statistical resolution}
Quadratic objectives are solved with auxiliary variables $z_i\ge |u_i-p_i|$ and simplex constraints. The implemented minimisation is scaled for numerical conditioning, but certificates are computed in the original objective units. Given a feasible iterate $x$ in the extended polytope $\mathcal P$, a separate linear programme gives the certificate
\begin{equation}
 C(x)=\max_{v\in\mathcal P}\nabla f(x)^\top(x-v),
\end{equation}
where $f$ is the convex negative reward. Convexity implies $f(x)-\min f\le C(x)$. This checks accuracy independently of the quadratic optimiser's success flag. Poorly certified iterates are retried at different scales and rejected if the certificate remains above $2\times10^{-7}$. The log comparators use a separate concave optimisation problem; the quadratic certificate is not claimed for them.

Realised monthly MV is the sample mean of net excess returns minus $\varphi$ times their sample variance. Excess Sharpe annualises the monthly mean-to-standard-deviation ratio by $\sqrt{12}$. These are evaluated on each strategy's complete own-state return path. For paired MV uncertainty, the same circular 12-month blocks are sampled for both strategies, with 1,000 repetitions and seed 8128. Shared blocks retain contemporaneous dependence between strategies and some serial dependence within each path. The intervals condition on the fitted strategies' realised paths: they do not repeat estimation, select hyperparameters or quantify every source of model uncertainty.

\section{Interpreting calibration across settings}
\subsection{Which comparison can establish improvement?}
A calibration rule is evaluated against a stated financial preference. In this study, improvement in alignment means a smaller $d_t^r$ or $g_t^r$ on the same estimated state. It does not mean that the calibrated strategy earns a higher return in the next month. The realised ranking can reverse because the target's estimated mean and covariance are imperfect forecasts. This distinction is especially relevant when the target itself uses a noisy expected-return estimate and a competing policy happens to reduce exposure to that estimate.

The primary distances put the scale of the comparison in context. Raw second moments give mean Euclidean distance 0.00843; unconstrained scalar calibration reduces it to 0.00540; reference-based risk-only calibration gives 0.00510. Joint calibration and the MVPI adaptation agree with the target to numerical precision. The log-minus-variance adaptation has distance 0.01286. These are distances between ten-dimensional weight vectors, not percentages of portfolio wealth lost. Fitted regret supplies the corresponding objective scale: the raw reward gap is only 0.0148 monthly basis points. A measurable allocation difference can therefore be economically small under the fitted preference.

The historical MV score answers a different question. Joint calibration is 1.315 monthly basis points below raw in the primary sample, while its interval spans both signs. That result does not support a claim that the calibrated policy is a superior investment strategy. Nor does it disprove the calibration identity: subsequent realised returns are absent from that identity. Minimum variance's better realised MV and Sharpe suggest that a strong practical implementation would need to address estimation quality alongside reward design.

\subsection{What the robustness grid tests}
The first block crosses three risk coefficients, $\varphi\in\{2,5,10\}$, with three turnover charges, $\lambda\in\{0,0.001,0.0025\}$, at the primary estimation window and shrinkage. The second crosses $L\in\{60,120,240\}$ with $h\in\{0,0.1,0.5\}$ at the primary risk and cost coefficients. Their shared primary setting is counted once, giving 17 distinct specifications. This is a union of two grids rather than the full cross-product of every choice. It probes economically relevant perturbations without pretending to span all market models.

The maximum setting-mean distance for joint calibration is $1.54\times10^{-7}$, compared with raw mean distances ranging from approximately 0.00157 to 0.03682. This supports the identity across the evaluated covariance estimates, constraints and cost levels. It does not supply 17 independent statistical replications: all settings reuse the same asset universe and overlapping return observations. The robustness runs compare the target, raw, scalar, risk-only, joint, equal-weight and minimum-variance policies. The MVPI and logarithmic comparators are evaluated in the primary study, rather than silently treated as rerun throughout this grid.

At zero transaction cost the risk-only and joint reward definitions coincide. With positive cost, the difference depends on the target expected return through $A=1+2\varphi y^*$. Constraints can make distinct objectives select the same boundary allocation, so exact agreement at an individual state need not establish equivalence of the rewards at all states. Conversely, a small discrepancy at the primary setting cannot be assumed to remain small when risk, cost or estimation changes.

\subsection{Optimisation effort and reward fidelity}
The finite-budget comparison assigns independent logits to each historical state. Softmax enforces strictly positive weights at finite logits, so exact boundary allocations can be approached without being attained. This helps explain why a high-accuracy constrained solver and a finite-budget softmax optimiser need not have equal residual errors even on the same objective. All three trained objectives use identical initialisation distributions, seeds and step budgets; the diagnostic compares numerical effort under that common parameterisation.

The own-reward gap measures whether the optimiser solved the reward it was given. The signed transfer term in the target decomposition instead measures how the approximate policy changes the target value relative to the exact reward optimum. It can have either sign. Adding own-reward suboptimality directly to target-unit design regret would mix different objective scales and would not reproduce the total target gap. At 5,000 steps, mean own-reward gaps remain around $4\times10^{-6}$, larger than the exact raw design gap of $1.48\times10^{-6}$. Improving the optimiser can therefore be the larger immediate numerical intervention, while calibration addresses the distinct discrepancy that persists at exact optimisation.

The complete robustness table and a reconstruction check appear in Appendices~\ref{app:robustness} and~\ref{app:verification}.

\section{Discussion}
The audit detects a specific preference change that aggregate performance can conceal. It also prevents overinterpreting its repair: minimum variance beats the fitted target in the primary historical sample, and the paired MV intervals do not establish a calibration advantage. Estimation risk remains after reward alignment has been resolved.

The procedure is suitable for a model-development validation record. Applying it in production would additionally require an agreed financial reference, point-in-time investable data, realistic execution costs and monitoring thresholds chosen for the mandate. Under return dependence, a terminal equilibrium also requires a continuation model; the rolling one-period reference used here does not supply one.

\bibliographystyle{ACM-Reference-Format}
\bibliography{references}
\appendix
\twocolumn[{
\section{Complete robustness results}\label{app:robustness}
Table~\ref{tab:fullrobust} reports all 17 saved specifications underlying Figure~\ref{fig:robust}. Each row averages Euclidean allocation distance over the same 432 historical months. These are repeated specifications on one dataset, not independent samples. The primary setting is $L=120$, $\varphi=5$, $\lambda=10$ bp and $h=0.1$. The three columns isolate an uncalibrated reward, calibration of risk alone, and joint calibration of risk and costs.

\begin{center}
\captionof{table}{Complete common-state alignment results. $L$ is the estimation window in months, $\lambda$ is the cost per unit of L1 turnover in basis points, and $h$ is diagonal covariance shrinkage. Distances are weight-vector norms, not returns. Bold identifies joint calibration (ours).}
\label{tab:fullrobust}
\begin{tabular}{rrrrrrr}\toprule
$L$ & $\varphi$ & $\lambda$ (bp) & $h$ & Raw distance & Risk-only distance & \textbf{Joint (ours)}\\\midrule
60 & 5 & 10 & 0 & 0.02541 & 0.01018 & \textbf{3.44e-08} \\
60 & 5 & 10 & 0.1 & 0.02391 & 0.00921 & \textbf{4.73e-08} \\
60 & 5 & 10 & 0.5 & 0.01747 & 0.00698 & \textbf{4.80e-08} \\
120 & 2 & 0 & 0.1 & 0.02032 & 0.00000 & \textbf{1.48e-07} \\
120 & 2 & 10 & 0.1 & 0.00697 & 0.00411 & \textbf{1.06e-08} \\
120 & 2 & 25 & 0.1 & 0.00376 & 0.00320 & \textbf{2.10e-09} \\
120 & 5 & 0 & 0.1 & 0.03664 & 0.00000 & \textbf{1.11e-07} \\
120 & 5 & 10 & 0 & 0.00874 & 0.00566 & \textbf{1.45e-08} \\
120 & 5 & 10 & 0.1 & 0.00843 & 0.00510 & \textbf{1.40e-08} \\
120 & 5 & 10 & 0.5 & 0.00822 & 0.00413 & \textbf{2.02e-08} \\
120 & 5 & 25 & 0.1 & 0.00324 & 0.00335 & \textbf{3.34e-09} \\
120 & 10 & 0 & 0.1 & 0.03682 & 0.00000 & \textbf{1.54e-07} \\
120 & 10 & 10 & 0.1 & 0.01296 & 0.00513 & \textbf{2.35e-08} \\
120 & 10 & 25 & 0.1 & 0.00394 & 0.00338 & \textbf{7.45e-09} \\
240 & 5 & 10 & 0 & 0.00271 & 0.00219 & \textbf{1.51e-09} \\
240 & 5 & 10 & 0.1 & 0.00237 & 0.00218 & \textbf{4.09e-09} \\
240 & 5 & 10 & 0.5 & 0.00157 & 0.00165 & \textbf{5.71e-09} \\
\bottomrule\end{tabular}
\end{center}
}]

\section{Independent reconstruction of the alignment check}\label{app:verification}
For a saved date, reconstruct the reference pre-trade holdings from the preceding reference allocation and gross asset returns. Recompute the lagged return moments using only observations before that date. Evaluate every saved audit allocation on this common state, including its own L1 distance from the reconstructed holdings. This gives the allocation distance and fitted regret without using the strategy's evolving own-state portfolio.

The calibration identity can also be checked away from the optimiser. For any feasible $u$, write $y=m^\top u^*$ and $A=1+2\varphi y$. The residual
\begin{equation}
 A Q(u)-F(u)+\varphi(m^\top u-y)^2-\varphi y^2
\end{equation}
should vanish. For the risk-only reward, add $2\varphi y\lambda\|u-p\|_1$ to this residual. These equalities test both coefficient scaling and the cost convention. They do not test whether estimated moments predict future returns accurately.

A separate saved verification sampled four simplex portfolios per primary historical state, using Dirichlet parameters equal to one and random seed 928. Across 1,728 sampled portfolios, the maximum absolute residual for each identity was below $9\times10^{-18}$. This numerical check supplements the algebraic proof; it is not a new market-performance experiment. The minimum primary calibration factor was 1.0208, satisfying the positive-factor assumption.

\end{document}
